% META: {"id": "TCOMPL-AIR-CR-010", "chapitre": "TCOMPL-CALCULS-AIRES", "type_objet": "cours", "capacites_codes": ["C4", "C5"], "sources_inspiration": [], "mode_creation": "ex_nihilo", "status": "generated"}

\section{Primitives d'une fonction}

% ============================================================
% STRATE 1 — Notion de primitive, formes usuelles
% ============================================================

\definition[1]{
Soit $f$ une fonction continue sur un intervalle $I$. Une \textbf{primitive} de $f$ sur $I$ est une fonction $F$ dérivable sur $I$ telle que $F'=f$.
}

\propriete[Primitives de formes usuelles]{
Si $u$ est une fonction dérivable sur $I$ :
\begin{itemize}
  \item une primitive de $2uu'$ est $u^2$ ;
  \item une primitive de $u'\mathrm{e}^u$ est $\mathrm{e}^u$ ;
  \item une primitive de $\dfrac{u'}{u}$ (avec $u>0$) est $\ln u$.
\end{itemize}
}

\exemple{
Une primitive de $f(x) = 2x\,\mathrm{e}^{x^2}$ sur $\mathbb{R}$ : on reconnaît la forme $u'\mathrm{e}^u$ avec $u(x)=x^2$, $u'(x)=2x$. Une primitive est donc $F(x) = \mathrm{e}^{x^2}$.
}

% BEGIN-VERIFY
% from sympy import symbols, exp, diff, simplify
% x = symbols('x')
% f = 2*x*exp(x**2)
% F = exp(x**2)
% assert simplify(diff(F, x) - f) == 0
% END-VERIFY

\margeAppui{
Reconnaître ces formes demande de repérer, dans l'expression à intégrer, une fonction $u$ et sa dérivée $u'$ qui apparaissent \textbf{ensemble} : c'est l'exercice principal du calcul de primitives au programme.
}

% ============================================================
% STRATE 2 — Deux primitives diferent d'une constante
% ============================================================

\theoreme[Primitives d'une même fonction]{
Si $F$ et $G$ sont deux primitives d'une même fonction $f$ continue sur un intervalle $I$, alors $F-G$ est \textbf{constante} sur $I$.
}

\demonstration{
Posons $H = F-G$. Alors $H' = F'-G' = f-f = 0$ sur $I$. Une fonction dont la dérivée est nulle sur un intervalle est constante sur cet intervalle (propriété admise, vue en première). Donc $H$ est constante, c'est-à-dire $F-G$ est constante sur $I$.
}

\exemple{
$F(x) = \dfrac{x^3}{3}$ et $G(x) = \dfrac{x^3}{3}+5$ sont deux primitives de $f(x)=x^2$ sur $\mathbb{R}$ : $F(x)-G(x) = -5$, bien constante.
}

% BEGIN-VERIFY
% from sympy import symbols, diff
% x = symbols('x')
% F = x**3/3
% G = x**3/3 + 5
% assert diff(F, x) == x**2
% assert diff(G, x) == x**2
% assert (F - G) == -5
% END-VERIFY

\erreurFrequente{
\textbf{Oublier la constante $+C$ dans une primitive générale.} Une fonction $f$ admet une \textbf{infinité} de primitives sur un intervalle, toutes de la forme $F+C$ où $F$ est une primitive particulière et $C$ une constante réelle quelconque : sans précision supplémentaire (par exemple une condition initiale), on ne peut pas déterminer \textbf{une seule} primitive.
}
